\section{Mixed products for a common measure}
\label{sec:common-measure-mixed-products}
\Cref{cor:bounded-multifactor-continuous-density} bounds a mixed Monge--Amp\`ere product from below when the pure Monge--Amp\`ere measures of the factors have continuous lower densities with respect to Lebesgue measure. For non-pluripolar products in \cref{chap:q09} we need the same inequality relative to an arbitrary positive Radon measure $\mu$, which may be singular with respect to Lebesgue measure. Instead of regularizing the potentials by convolution, as in the proof of that corollary, we restrict each pure Monge--Amp\`ere measure to a compact set and replace it by its cube averages, as in \cref{prop:canonical-grid-energy-control}. The averaged measures have piecewise constant densities, so the case of Lebesgue measure applies on each cube, and H\"older's inequality compares the resulting bound with the cube averages of the expected lower bound. The convergence results of \cref{sec:local-energy-moments} then allow us to pass to the limit.

\begin{theorem}\label{thm:bounded-multifactor-common-measure}
Let $\Omega\subset\mathbb C^n$ be open, $n\geq1$, and let $u_1,\ldots,u_n$ be locally bounded plurisubharmonic functions on $\Omega$. Let $\mu$ be a positive Radon measure on $\Omega$, and let $f_1,\ldots,f_n$ be nonnegative Borel functions on $\Omega$. If
\[
 (\ddc u_i)^n\geq f_i\mu\qquad(1\leq i\leq n),
\]
then
\begin{equation}\label{eq:bounded-multifactor-common-measure}
 \ddc u_1\wedge\cdots\wedge\ddc u_n
                  \geq\left(\prod_{i=1}^nf_i^{1/n}\right)\mu.
\end{equation}
For $n=1$ the assertion is the hypothesis.
\end{theorem}
\begin{proof}
Let $n\geq2$. We use the generalized H\"older inequality
\begin{equation}\label{eq:generalized-holder-geometric-mean}
 \int_A\prod_{i=1}^ng_i^{1/n}\,\mathrm d\nu\leq\prod_{i=1}^n\Bigl(\int_Ag_i\,\mathrm d\nu\Bigr)^{1/n}
\end{equation}
for a positive measure $\nu$, a measurable set $A$ and nonnegative measurable functions $g_i$. If one of the integrals $c_i=\int_Ag_i\,\mathrm d\nu$ vanishes, then $g_i=0$ $\nu$-almost everywhere on $A$ and both sides vanish, with the convention $0\cdot\infty=0$; if none vanishes and one is infinite, the right side is infinite. Otherwise, integrating the arithmetic--geometric mean inequality $\prod_i(g_i/c_i)^{1/n}\leq\frac1n\sum_ig_i/c_i$ over $A$ gives $\int_A\prod_i(g_i/c_i)^{1/n}\,\mathrm d\nu\leq1$, which is \eqref{eq:generalized-holder-geometric-mean}.\index{H\"older inequality!generalized} For every compact $K\subset\Omega$ we have $\int_Kf_i\,\mathrm d\mu\leq(\ddc u_i)^n(K)<\infty$ by hypothesis, so \eqref{eq:generalized-holder-geometric-mean} shows that the right side of \eqref{eq:bounded-multifactor-common-measure} is a positive Radon measure. We first prove the theorem when $\Omega$ is a ball $B=B(a_0,R)$ and all $u_i$ belong to $\mathcal E_0(B)$ (Steps 1 and 2), and then reduce the general case to this one by modifying the potentials near the boundary of a ball (Step 3). Put $\rho(z)=|z-a_0|^2-R^2$.

\emph{Step 1. Exhaustion by compactly supported measures.}
Choose closed balls $K_l\subset B$ concentric with $B$, increasing with $l$ and exhausting $B$, and set
\[
 \alpha_{i,l}=\mathbf1_{K_l}(\ddc u_i)^n.
\]
These measures are finite since $u_i\in\mathcal E_0(B)$, they have compact support in $B$, and they are dominated by the Monge--Amp\`ere measure of the bounded plurisubharmonic function $u_i$. By \cref{thm:compact-bounded-subsolution-ball-dirichlet}, the equation $(\ddc A)^n=\alpha_{i,l}$ has a bounded solution $A_{i,l}\in\mathcal E_0(B)$ with zero boundary values, and $A_{i,l}$ is the largest bounded plurisubharmonic function $v$ on $B$ with $(\ddc v)^n\geq\alpha_{i,l}$ and upper boundary limits at most zero. The function $u_i$ has these properties for every $l$, and so has $A_{i,l+1}$, because $\alpha_{i,l+1}\geq\alpha_{i,l}$. Therefore
\[
 u_i\leq A_{i,l+1}\leq A_{i,l}\leq0.
\]

We claim that $A_{i,l}\downarrow u_i$ as $l\to\infty$. The decreasing limit $A_{i,\infty}$ satisfies $u_i\leq A_{i,\infty}\leq0$, so it is a bounded plurisubharmonic function (\cref{thm:psh-decreasing-nets}) with zero boundary values. By \cref{thm:bt-decreasing-continuity}, $(\ddc A_{i,\infty})^n$ is the weak limit of $\alpha_{i,l}$, that is,
\[
 (\ddc A_{i,\infty})^n=(\ddc u_i)^n,
\]
and this measure has finite total mass. Let $\varepsilon>0$ and $S=\{u_i<A_{i,\infty}+\varepsilon\rho\}$. Since $u_i-(A_{i,\infty}+\varepsilon\rho)$ tends to zero at $\partial B$, the comparison theorem (\cref{thm:bt-comparison}) applies. Expanding $(\ddc(A_{i,\infty}+\varepsilon\rho))^n$ and dropping the mixed terms, which are positive, we obtain
\[
 \int_S(\ddc u_i)^n+\varepsilon^n\int_S(\ddc\rho)^n
 \leq\int_S\bigl(\ddc(A_{i,\infty}+\varepsilon\rho)\bigr)^n
 \leq\int_S(\ddc u_i)^n.
\]
The integrals of $(\ddc u_i)^n$ are finite, so $S$ is Lebesgue null. Hence $u_i\geq A_{i,\infty}+\varepsilon\rho$ almost everywhere, and therefore everywhere on $B$, by the uniqueness assertion of \cref{thm:psh-distribution-representative} applied to $\max\{u_i,A_{i,\infty}+\varepsilon\rho\}$ and $u_i$. Letting $\varepsilon\downarrow0$ gives $A_{i,\infty}=u_i$, as claimed.

\emph{Step 2. A common grid.}
Fix $l$. Choose a sequence of partitions of $\mathbb C^n=\mathbb R^{2n}$ into half-open cubes whose side lengths decrease to zero, the same for all $i$. Let $\alpha_{i,l,j}$ be the corresponding cube averages of $\alpha_{i,l}$, as in \cref{prop:canonical-grid-energy-control}, and let $A_{i,l,j}\in\mathcal E_0(B)$ be the corresponding canonical approximants of $\alpha_{i,l}$, that is, the zero-boundary solutions of $(\ddc A_{i,l,j})^n=\alpha_{i,l,j}$; they are defined for all large $j$. Discarding finitely many indices $j$, we may assume that they are defined for all $j$ and all $i$, and that every cube meeting $K_l$ has its closure in $B$. By \cref{thm:compact-bounded-subsolution-ball-dirichlet}, $A_{i,l,j}\to A_{i,l}$ locally in capacity as $j\to\infty$. By \cref{prop:canonical-grid-energy-control}, applied with the bounded subsolution $u_i$, the energies $e_1(A_{i,l,j})$ are bounded uniformly in $j$; since there are only finitely many indices $i$, they have a common bound. (If $\alpha_{i,l}=0$, all these solutions vanish.) Therefore \cref{cor:mixed-products-energy-capacity-continuity} gives
\begin{equation}\label{eq:common-measure-grid-mixed-limit}
 \bigwedge_{i=1}^n\ddc A_{i,l,j}
                      \longrightarrow\bigwedge_{i=1}^n\ddc A_{i,l}
\end{equation}
weakly on $B$.

Let $\lambda_l=\mathbf1_{K_l}\prod_if_i^{1/n}\mu$, a finite measure, and let $\lambda_{l,j}$ be its cube averages. For every cube $Q$, the hypothesis gives $\alpha_{i,l}(Q)\geq\int_{Q\cap K_l}f_i\,\mathrm d\mu$, and \eqref{eq:generalized-holder-geometric-mean} then gives
\begin{equation}\label{eq:common-measure-grid-holder}
 \prod_{i=1}^n\left(\frac{\alpha_{i,l}(Q)}{|Q|}\right)^{1/n}
 \geq\frac1{|Q|}\int_{Q\cap K_l}\prod_{i=1}^nf_i^{1/n}\,\mathrm d\mu
 =\frac{\lambda_l(Q)}{|Q|}.
\end{equation}
Let $Q$ be a cube with $\lambda_l(Q)>0$; it meets $K_l$, so its interior lies in $B$. There the measures $(\ddc A_{i,l,j})^n=\alpha_{i,l,j}$ have the constant densities $\alpha_{i,l}(Q)/|Q|$ with respect to Lebesgue measure $\mathrm dV$. Recall that $\eta^n=c_n\,\mathrm dV$ with $\eta=\ddc|z|^2$ and $c_n=n!\pi^{-n}$, by \eqref{eq:bt-determinant-normalization}. Applying \cref{cor:bounded-multifactor-continuous-density} on the open set $\Int Q$ to the constant densities $c_n^{-1}\alpha_{i,l}(Q)/|Q|$ with respect to $\eta^n$, and using \eqref{eq:common-measure-grid-holder}, we obtain on $\Int Q$
\[
 \bigwedge_{i=1}^n\ddc A_{i,l,j}
 \geq\prod_{i=1}^n\left(\frac{\alpha_{i,l}(Q)}{c_n|Q|}\right)^{1/n}\eta^n
 =\prod_{i=1}^n\left(\frac{\alpha_{i,l}(Q)}{|Q|}\right)^{1/n}\mathrm dV
 \geq\lambda_{l,j}.
\]
The factor $c_n$ cancels because the $n$ exponents $1/n$ sum to one. The measure $\lambda_{l,j}$ vanishes on the cubes with $\lambda_l(Q)=0$ and does not charge the cube boundaries, which are Lebesgue null. Since the mixed product is positive, we conclude that
\[
 \bigwedge_{i=1}^n\ddc A_{i,l,j}\geq\lambda_{l,j}\quad\text{on }B.
\]

As $j\to\infty$, $\lambda_{l,j}\to\lambda_l$ weakly. Indeed, only finitely many cubes have positive $\lambda_l$-mass, and for $\chi\in C_c(B)$, replacing the values of $\chi$ by its cube averages changes $\int\chi\,\mathrm d\lambda_l$ by at most $\lambda_l(B)$ times the modulus of continuity of $\chi$ at the cube diameter. Letting $j\to\infty$ and using \eqref{eq:common-measure-grid-mixed-limit}, we obtain $\bigwedge_i\ddc A_{i,l}\geq\lambda_l$. Finally let $l\to\infty$. By Step 1, $A_{i,l}\downarrow u_i$ with $u_i\leq A_{i,l}\leq0$, so \cref{thm:bt-decreasing-continuity} gives $\bigwedge_i\ddc A_{i,l}\to\bigwedge_i\ddc u_i$ weakly, while $\lambda_l$ increases to $\prod_if_i^{1/n}\mu$. Testing with nonnegative functions in $C_c(B)$ and using monotone convergence on the right, we obtain \eqref{eq:bounded-multifactor-common-measure} for potentials in $\mathcal E_0(B)$.

\emph{Step 3. Localization.}
Now let $u_1,\ldots,u_n$ be locally bounded on $\Omega$. Choose a ball $B=B(a_0,R)$ with $\overline B\subset\Omega$ and a compact set $K\subset B$. Since adding constants does not change the products, we may assume that $u_i\leq-1$ on $B$. With $\rho(z)=|z-a_0|^2-R^2$ and constants $C_i>0$ so large that $C_i\rho<\inf_Bu_i$ on a neighborhood of $K$, put
\[
 \widetilde u_i=\max\{u_i,C_i\rho\}.
\]
Then $\widetilde u_i=u_i$ on a neighborhood of $K$, and $\widetilde u_i=C_i\rho$ near $\partial B$, because $C_i\rho\to0$ at $\partial B$ while $u_i\leq-1$. Thus $\widetilde u_i\in\mathcal E_0(B)$: its Monge--Amp\`ere measure is smooth near $\partial B$ and locally finite on the rest of $B$. By locality of the products on the open neighborhood of $K$ where $\widetilde u_i=u_i$, we have $(\ddc\widetilde u_i)^n\geq f_i\mathbf1_K\mu$ on $B$. The case already proved, applied to the measure $\mathbf1_K\mu$, gives $\bigwedge_i\ddc\widetilde u_i\geq\prod_if_i^{1/n}\mathbf1_K\mu$, and on a neighborhood of $K$ the left side equals $\bigwedge_i\ddc u_i$ by locality. Thus $P=\bigwedge_i\ddc u_i$ and $\lambda=\prod_if_i^{1/n}\mu$ satisfy $P(E)\geq\lambda(E)$ for every Borel set $E\subset K$. The interiors of such compact sets $K$ cover $\Omega$, and since $\Omega$ has a countable basis, countably many of them, say $\Int K_1,\Int K_2,\ldots$, already cover $\Omega$. A Borel set $E\subset\Omega$ is the disjoint union of the Borel sets $E_m=(E\cap K_m)\setminus\bigcup_{m'<m}K_{m'}$, and therefore $P(E)=\sum_mP(E_m)\geq\sum_m\lambda(E_m)=\lambda(E)$.
\end{proof}
\index{Monge--Amp\`ere measure!multifactor common-measure inequality}

\begin{corollary}\label{cor:bounded-common-measure-mixed}
Let $u,v$ be locally bounded plurisubharmonic functions on an open set $\Omega\subset\mathbb C^n$, $n\geq1$, and let $\mu$ be a positive Radon measure on $\Omega$. If $f,g$ are nonnegative Borel functions on $\Omega$ and
\[
 (\ddc u)^n\geq f\mu,\qquad(\ddc v)^n\geq g\mu,
\]
then, for each integer $1\leq k\leq n-1$,
\begin{equation}\label{eq:bounded-common-measure-mixed}
 (\ddc u)^k\wedge(\ddc v)^{n-k}
                       \geq f^{k/n}g^{(n-k)/n}\mu.
\end{equation}
For $k=0$ and $k=n$ the corresponding inequalities are the hypotheses; for $n=1$ only these occur.
\end{corollary}
\begin{proof}
Apply \cref{thm:bounded-multifactor-common-measure} to $k$ copies of $u$ with density $f$ and $n-k$ copies of $v$ with density $g$.
\end{proof}
\index{Monge--Amp\`ere measure!common-measure mixed inequality}

\begin{corollary}\label{cor:bounded-common-measure-sum}
Under the hypotheses of \cref{cor:bounded-common-measure-mixed},
\[
 (\ddc(u+v))^n\geq(f^{1/n}+g^{1/n})^n\mu.
\]
\end{corollary}
\begin{proof}
By additivity and symmetry of the products, $(\ddc(u+v))^n=\sum_{k=0}^n\binom nk(\ddc u)^k\wedge(\ddc v)^{n-k}$. Bound the terms with $1\leq k\leq n-1$ by \cref{cor:bounded-common-measure-mixed} and the terms with $k=0,n$ by the hypotheses; the binomial formula then gives the assertion. The density on the right is locally $\mu$-integrable, since $f^{k/n}g^{(n-k)/n}\leq f+g$ by the weighted arithmetic--geometric mean inequality, and $f$ and $g$ are locally $\mu$-integrable because $f\mu\leq(\ddc u)^n$ and $g\mu\leq(\ddc v)^n$.
\end{proof}
\index{Monge--Amp\`ere measure!common-measure sum inequality}
